% CX_GRH: standalone section fragment. No preamble.
% Bibliography keys are supplied in koiran_bibliography.tex.
% Checked 2026-11-20 UTC by /root/koiran_gct, GPT-6 Astra Pro,
% inherited effort (setting not exposed). Independent full-proof checks:
% /root/eth_consequences and /root, with qualifications incorporated.

\section{Koiran and Noether normalization: the exact analytic gate}

The analytic hypothesis needed by Koiran's proof is RH for the Dedekind
zeta functions of bounded solution fields. These fields can be arbitrary
number fields. They need not be Galois over $\mathbb Q$, and passing to a
Galois closure is unnecessary. Masullo's all-number-field Hecke-GRH claim,
if accepted as $H_M$, supplies this hypothesis by its trivial-character
case. The stipulated Alp\"oge Dirichlet-GRH theorem $H_A$ supplies the
corresponding input for abelian extensions of $\mathbb Q$, but the checked
factorization argument does not extend this to all solution fields.

Write
\[
 H_D:\quad \text{RH holds for $\zeta_E$ for every number field $E$.}
\]
Throughout this section, $H_M$ and $H_A$ have the precise meanings fixed
in the premise register. In particular, $H_M$ remains a conditional
claim. Our conclusions under $H_M$ are transfers of established
complexity theorems, together with the encoding and scope refinements
proved below. They are not independent unconditional proofs of those
theorems.

\subsection{Source conventions, numerical locators, and representations}

Koiran's published paper and extended report require separate locators.
The journal version \cite[\S2, Theorem~2, pp.~276--277]{K96} displays
$\mathrm{RP}^{\mathrm{NP}}$, followed by containment in
$\Pi_2^{\mathrm P}$. The extended DIMACS report explicitly displays
$\mathrm{AM}\subseteq\Pi_2^{\mathrm P}$ in \S2, Theorem~2,
printed pp.~5--6 \cite{K96TR}. Its protocol uses public randomness before
the nondeterministic collision witnesses. We use this report for the
requested AM formulation. Koiran's 8 August 1997 erratum, first page,
second bullet, corrects the lower bound in Theorem~8 and states that
Theorem~1 and the remainder of the paper are unaffected
\cite{K97Err}. All transfers here use the corrected version.

The input in Koiran's \S1.1 is a finite list of integer polynomials in
the Turing model, with sparse encoding. Theorems~4 and~7, Lemma~3, and
Theorem~8 supply the bounded algebraic witness and reduction modulo
primes. Theorem~9 and Corollary~1 in \S5, journal p.~282 and DIMACS
printed p.~12, isolate the effective prime-ideal theorem. The relevant
analytic result is Lagarias--Odlyzko \cite[Theorem~1.1]{LO77}, whose
hypothesis is RH for the Dedekind zeta function of the top field.

For GCT~V \cite{GCT5}, the attached \texttt{main.tex} labels and the
numerical statements of arXiv:1209.5993v5 agree as follows. Page numbers
in this table refer to that arXiv PDF, not the journal pagination.

\begin{center}
\small
\begin{tabular}{@{}p{.14\linewidth}p{.25\linewidth}p{.43\linewidth}r@{}}
\hline
Statement & Attached source label & Use & Page \\
\hline
Theorem 2.8 & \texttt{tstrongexp} & Strengthened black-box PIT, EXPH under GRH & 24--25 \\
Proposition 2.9 & \texttt{pstandard} & Standard black-box PIT, PH under GRH & 24 \\
Theorem 2.10 & \texttt{tkoiran} & Circuit and homogeneous HN & 24 \\
Lemma 3.1 & \texttt{lnnl0} & Geometric normalization & 26 \\
Lemma 3.2 & \texttt{lnoether20} & Graded integrality criterion & 26 \\
Theorem 3.4 & \texttt{tnoetherph} & Equation-given h.s.o.p., functional PH & 27--28 \\
Theorem 4.1 & \texttt{tnnldetexpspace} & Determinant h.s.o.p., extra equation hypothesis & 30 \\
Theorem 4.10 & \texttt{tssopverifydet} & Determinant s.s.o.p., construction/verification & 34 \\
Definition 5.1 & \texttt{dexpvariety} & Explicit families and succinct circuits & 35--36 \\
Theorem 5.5 & \texttt{tnnlexplicithsop} & Explicit h.s.o.p., extra equation hypothesis & 39 \\
Theorem 5.11(a) & \texttt{tnnlexpfull} & Border hitting sets give strict separating systems & 43 \\
Theorem 5.12 & \texttt{tssopexplicit} & Explicit strict s.s.o.p., construction/verification & 44 \\
Theorem 7.3 & \texttt{thsop} & Simultaneous-conjugation invariant ring & 50 \\
\hline
\end{tabular}
\end{center}

An h.s.o.p. has exactly $\dim R$ algebraically independent homogeneous
parameters and makes $R$ finite over their polynomial subring. A small
system of parameters may have polynomially many elements, with no
optimal-cardinality requirement. The succinct small systems in GCT~V
are specified through small circuits and are a different output regime
from explicitly listing exponentially many ambient coordinates.
``Functional PH'' means a construction with polynomial-time oracle
access to a fixed level of PH. It is not a decision-language assertion
and does not mean ordinary polynomial-time construction.

The standard specification preceding GCT~V Theorem~3.4 is a homogeneous
affine cone defined by homogeneous polynomial equations represented by
arithmetic circuits with rational constants. The characteristic-zero
statements here use this effective Turing representation. We do not
assert Turing bounds for arbitrarily presented transcendental
coefficients. The stronger refinement in CX-K05 explicitly uses
standard grading. GCT~V Definition~5.1 also allows differing positive
weights; the cited mixed-weight theorem is kept separate in CX-K06.

More precisely, Definition~5.1 presents
$W_n=\overline{\psi_n(K^{r_n})}\subseteq K^{m_n}$ with
$r_n=\operatorname{poly}(n)$, $m_n=n^{\Omega(1)}$, and
$\log m_n=\operatorname{poly}(n)$. Thus
$\dim W_n\le r_n$, although its ambient dimension may be exponential.
Writing $\psi_n=(f_1,\ldots,f_{m_n})$, the succinct input is a
uniformly constructible rational circuit, of polynomial total bit
length, for $F_n(v,x)=\sum_j f_j(v)g_j(x)$. The $g_j$ are linearly
independent; the $f_j$ and $g_j$ are homogeneous of polynomial degree,
and $F_n$ has polynomial total degree. Definition~5.6 measures a short output by circuits for
$\psi_n^*(s)$, including their rational constants' bit lengths.
Definition~5.7 calls an element strict when it is a specialization
of a fixed homogeneous-in-$v$ component of $F_n$ at a polynomial-bit
integer point $x=b$. Hence a strict output is encoded by polynomially
many such points and component indices, without listing $m_n$
ambient coefficients. The dimension bound alone does not require the
small system's cardinality to equal the dimension.

\subsection{CX-K01: the root-field gate without a Galois closure}
\label{cx:k01}

\paragraph{Statement.}
In Koiran's proof, the analytic premise can be localized inputwise to
RH for $\zeta_E$, where $E$ is the field of a bounded algebraic solution
chosen using Koiran's Theorems~4 and~7. The uniform family of possible
$E$ is unrestricted. Neither Artin holomorphy nor passage to a splitting
field is required. If one uses $\operatorname{Li}(x)$ rather than
Koiran's comparison with the rational prime-counting function, no
separate RH-for-$\mathbb Q$ premise is needed at this step.

\begin{proof}
Let $S$ be a satisfiable integer polynomial system of encoding length
$N$. Koiran's Theorems~7 and~4 give a solution
$a=(a_1,\ldots,a_t)$, a primitive element $\theta$ for
$E=\mathbb Q(a_1,\ldots,a_t)$, an integer minimal polynomial $R$ of
$\theta$, and integer denominators in expressions for the $a_i$.
The degree, logarithmic discriminant bounds, and denominator logarithms
are all at most $2^{q(N)}$ for a fixed polynomial $q$.

We can make the primitive element integral without losing that scale.
If $R$ has leading coefficient $c$ and degree $d$, then
$c^{d-1}R(T/c)$ is a monic integral polynomial for
$\beta=c\theta$. Thus $\mathbb Z[\beta]\subseteq\mathcal O_E$
has finite index, the field discriminant satisfies the usual
polynomial-discriminant/index relation, and all $a_i$ lie in
$\mathcal O_E[1/B]$ for an integer $B\ne0$ with singly exponentially
bounded logarithm. Put $J=[\mathcal O_E:\mathbb Z[\beta]]$ and
$D_E=\operatorname{disc}(E)$. The logarithm of $J$ has the same bound.

Apply Lagarias--Odlyzko Theorem~1.1 to the trivial extension $E/E$.
It is a Galois extension, regardless of whether $E/\mathbb Q$ is
normal. Under RH for $\zeta_E$, the number $\pi_E(x)$ of prime ideals
of $\mathcal O_E$ of norm at most $x$ satisfies
\begin{equation}
 \pi_E(x)=\operatorname{Li}(x)+
 O\!\left(\sqrt{x}\log(|D_E|x^d)+\log |D_E|\right),
 \qquad d=[E:\mathbb Q],
 \label{eq:k01primeideal}
\end{equation}
with an effective absolute implied constant. This is an invocation of
the cited analytic theorem, not a new proof of that theorem.

A prime ideal of residue degree at least two and norm at most $x$
lies above a rational prime at most $\sqrt{x}$. At most $d$ prime
ideals lie above any one rational prime. Hence these higher-degree
prime ideals contribute at most $d\pi_{\mathbb Q}(\sqrt{x})\le
d\sqrt{x}$. Discard also the prime ideals above divisors of $BJ$;
there are at most $d\omega(BJ)$, where $\omega$ counts distinct prime
divisors. At any remaining residue-degree-one prime ideal
$\mathfrak p$, reduction maps the solution coordinates into
$\mathcal O_E/\mathfrak p=\mathbb F_p$. The equations of $S$ remain
valid. Since each rational prime has at most $d$ prime ideals above it,
the number of rational primes $p\le x$ for which $S$ is soluble is at
least
\begin{equation}
 \frac{1}{d}\left[
 \operatorname{Li}(x)
 -O\!\left(\sqrt{x}\log(|D_E|x^d)+d\sqrt{x}+\log|D_E|\right)
 -d\omega(BJ)\right].
 \label{eq:k01usefulprimes}
\end{equation}

All quantities outside $x$ and $\log x$ in this expression are bounded
by $2^{\operatorname{poly}(N)}$. Choose
$x=2^{P(N)}$ for a sufficiently large fixed polynomial $P$.
The main term in \eqref{eq:k01usefulprimes} then dominates its errors
and the unconditional inconsistent-system bound in Koiran's
Theorem~5 by the gap required in his Theorem~1. Its primes still have
polynomial bit length. Koiran's unchanged public-coin protocol in
DIMACS Theorem~2 therefore applies. This identifies the exact analytic
gate and shows that no Galois closure is needed. In particular, the
potential factorial degree growth of a splitting field is avoided.

Finally, every number field $E$ has a primitive element $\theta$.
Clear denominators in its irreducible polynomial $R$. All roots of
$R$ generate conjugate copies of $E$. Thus the single equation
$R(x)=0$ realizes this isomorphism type among the possible solution
fields. The input family is therefore not confined to abelian fields.
This observation concerns the fields occurring in the checked proof;
it does not establish that the complexity conclusion logically
requires $H_D$.
\end{proof}

The weaker unconditional prime-ideal estimate in
Lagarias--Odlyzko Theorem~1.3 starts in the range
$x\ge\exp(10d(\log|D_E|)^2)$. With singly exponential $d$ and
$\log|D_E|$, that range need not preserve polynomial bit length of
the primes. It therefore does not substitute into this proof with the
same complexity bound. This is a limitation of the checked route,
not an independence claim about Dirichlet GRH and $\mathrm{HN}\in
\mathrm{AM}$.

\subsection{CX-K02: conditional HN for circuits and nonzero homogeneous points}
\label{cx:k02}

\paragraph{Statement.}
Under $H_M$, feasibility over $\mathbb C$ of polynomial systems given
by arithmetic circuits with rational constants belongs to
\[
 \mathrm{AM}\subseteq\Pi_2^{\mathrm P}.
\]
The same bound holds for existence of a nonzero solution of a
homogeneous such system. The complementary infeasibility predicates
belong to $\Sigma_2^{\mathrm P}$.

\begin{proof}
For every number field $E$, the Hecke $L$-function of the trivial
character is $\zeta_E$. Consequently $H_M\Rightarrow H_D$.
Apply the corrected Koiran DIMACS Theorem~2, or CX-K01 followed by its
unchanged hashing protocol. This gives the asserted AM bound for
Koiran's sparse input representation.

For a circuit representation, introduce a variable $y_g$ for each
gate. A sum gate and a product gate give, respectively,
\[
 y_g-y_h-y_j=0,\qquad y_g-y_hy_j=0.
\]
A rational constant $a/b$ gives $by_g-a=0$, with $b\ne0$.
Wire the input gates to the original variables, and set each output
gate variable to zero. The result is a polynomial-size sparse system
of degree at most two and polynomial total coefficient bit length.
Every original solution has exactly one extension to gate values.
Conversely, induction over the circuit forces all gate variables to
equal the intended evaluations. This is a deterministic
polynomial-time many-one reduction, which preserves AM membership.

For a homogeneous system in $x_1,\ldots,x_t$, adjoin variables
$u_1,\ldots,u_t$ and the equation
\[
 \sum_{i=1}^t u_i x_i=1.
\]
A nonzero solution admits such a $u$ by inverting a nonzero coordinate,
and a solution to this enlarged system necessarily has $x\ne0$.
This again is a polynomial-size circuit reduction. The standard
containment used in Koiran is $\mathrm{AM}\subseteq\Pi_2^{\mathrm P}$;
complementing gives $\Sigma_2^{\mathrm P}$. No containment
$\mathrm{AM}\subseteq\Sigma_2^{\mathrm P}$ is asserted.
\end{proof}

\subsection{CX-K03: the bounded abelian-witness consequence of the stipulated Dirichlet theorem}
\label{cx:k03}

\paragraph{Statement.}
Fix an explicit polynomial $P$. Consider the promise problem whose
no-instances are all inconsistent rational circuit systems and whose
yes-instances have a solution as in CX-K01 with a finite abelian
Galois field $E/\mathbb Q$ and degree, logarithmic defining-polynomial
discriminant, and denominator logarithms at most $2^{P(N)}$.
Under $H_A$, this promise problem is in AM. Its verifier need not
compute $E$, certify abelianness, or recognize the promise.

\begin{proof}
For finite abelian $E/\mathbb Q$, the Dedekind zeta function factors
as the product of primitive Dirichlet $L$-functions attached to
characters of $\operatorname{Gal}(E/\mathbb Q)$, with the conductor
convention giving the correct ramified Euler factors. The trivial
factor is $\zeta_{\mathbb Q}$. The stipulated $H_A$ therefore puts
all nontrivial zeros of $\zeta_E$ on the critical line.

On every promised yes-instance, CX-K01 now gives the useful-prime
bound with a uniform polynomial exponent. On every no-instance,
Koiran's Theorem~5 supplies the unconditional bound. The same gap
and public-coin protocol used in Koiran's DIMACS Theorem~2 apply.
The arithmetic promise proves completeness of the analytic estimate;
the protocol never checks it. Soundness on all inconsistent systems
is unchanged.
\end{proof}

The bound on the abelian witness is essential to the statement proved
here. Koiran's existence of a bounded solution does not say that a
solution satisfying an additional arithmetic property can be chosen
with that same bound. We therefore do not replace the quantitative
promise by the bare assertion that some abelian-valued solution exists.

\subsection{CX-K04: a counterexample to replacing all solution fields by abelian fields}
\label{cx:k04}

\paragraph{Statement.}
The satisfiable equation $x^3-2=0$ has no solution in any abelian
extension of $\mathbb Q$. Hence forcing every witness into an abelian
field is not an available general reduction.

\begin{proof}
Eisenstein's criterion at $2$ makes $x^3-2$ irreducible. Its real root
generates a degree-three subfield of $\mathbb R$ that omits the two
nonreal conjugates, so that root field is not normal over $\mathbb Q$.
The other two root fields are conjugate to it and are likewise
nonnormal. Every intermediate field of a finite abelian Galois
extension is Galois over $\mathbb Q$, since every subgroup of an
abelian group is normal. None of these root fields can be an
intermediate field. A root in an infinite abelian extension would
already lie in a finite abelian subextension, so that possibility
also fails.
\end{proof}

Only the abelian-witness shortcut is refuted. This counterexample does
not refute the possibility of a different proof of
$\mathrm{HN}\in\mathrm{AM}$ from $H_A$ alone.

\subsection{CX-K05: the functional PH level for equation-given h.s.o.p.'s}
\label{cx:k05}

\paragraph{Statement.}
Assume $H_M$. Let
$R=\mathbb Q[z_1,\ldots,z_t]/I$ be a nonzero standard positively
graded algebra, with $I$ specified by homogeneous rational-constant
arithmetic-circuit equations. Let $N$ be the total encoding length,
including an explicit list of the ambient coordinates, so $t\le N$.
An h.s.o.p. of rational linear forms and $\dim R$ can be constructed
in
\[
 \mathrm{FP}^{\Sigma_2^{\mathrm P}}
\]
with polynomial output bit length. These forms remain an h.s.o.p.
after extension of scalars to any characteristic-zero field.
Unconditionally the construction has a polynomial-space bound.

\begin{proof}
The bounded-coefficient part of GCT~V Theorem~3.4's proof gives a
computable polynomial $b(N)$ such that, for $d=\dim R$, there are
$d$ integral linear forms, each coefficient having at most $b(N)$
bits, whose common zero on the cone is only its vertex. The relevant
size argument in that proof bounds the equation degrees by $2^N$,
passes to a common homogeneous degree at most $2^{N^2}$, and uses a
nonzero resultant of degree with polynomial logarithm. In particular,
only a polynomial bit bound is needed for the coefficient grid;
enumerating all monomials is unnecessary. Nilpotents do not change
the set of geometric zeros or the dimension, so the same existence
bound applies to the algebra stated here.

For $L=(L_1,\ldots,L_s)$ let $\operatorname{Good}(I,L)$ mean that
the equations generating $I$, together with $L_1=\cdots=L_s=0$,
have no nonzero geometric solution. By CX-K02 this is a
$\Sigma_2^{\mathrm P}$ predicate. Fix a canonical binary encoding
of the bounded coefficient lists, padded to polynomial length. For
any prescribed prefix, the predicate
\[
 \exists L\;[\text{$L$ extends the prefix and }
                    \operatorname{Good}(I,L)]
\]
also belongs to $\Sigma_2^{\mathrm P}$: the bounded existential
guess of $L$ merges with the leading existential block. In explicit
quantifier form, if Good is $\exists u\,\forall v\,T(I,L,u,v)$
for polynomial-time $T$, the prefix predicate is
$\exists (L,u)\,\forall v\,[\operatorname{Prefix}(L)\land
T(I,L,u,v)]$.

Ask this predicate for $s=0,1,\ldots,t$. The least successful $s$
equals $d$. Existence at $d$ follows from the bounded-coefficient
argument. For $s<d$, the generalized principal ideal theorem gives
\[
 \dim R/(L_1,\ldots,L_s)\ge d-s>0,
\]
so a homogeneous cone cannot be supported only at the origin.
For $d=0$, the empty list works.

Once $s=d$ is known, recover a successful list one bit at a time.
Extend the current prefix by zero and ask if a successful completion
exists. Keep zero when the answer is yes, and otherwise keep one.
The invariant that the prefix has a completion holds initially and
is preserved by every step. There are polynomially many coefficient
bits and polynomially many dimension queries. This proves the
$\mathrm{FP}^{\Sigma_2^{\mathrm P}}$ bound, rather than merely a
language membership in PH.

We verify that the searched geometric condition is precisely the
required ring condition. Good makes $R/(L)$ finite-dimensional over
$\mathbb Q$: after extending to an algebraic closure, its spectrum is
supported at the origin, so its homogeneous irrelevant ideal is
nilpotent. Choose homogeneous representatives of a finite vector-space
basis of $R/(L)$. They generate $R$ over
$A=\mathbb Q[L_1,\ldots,L_d]$. Indeed, for any homogeneous element,
subtract a representative of its class modulo $(L)$. The difference
is a sum of terms $L_i r_i$, with each $r_i$ of strictly lower degree;
induction on degree proves the assertion. Thus $R$ is a finite
$A$-module and is integral over $A$.

Conversely, if $R$ is finite over $A$, its quotient by the images of
the $L_i$ is finite-dimensional over $\mathbb Q$, so its homogeneous
cone has no nonzero geometric point. Finally, integral extensions
preserve dimension, so $\dim A=d$. A nonzero kernel in
$\mathbb Q[T_1,\ldots,T_d]\to A$ would make $\dim A<d$.
Thus the parameters are algebraically independent and constitute an
h.s.o.p. This reasoning permits nilpotents and several irreducible
components.

All constructed coefficients are rational. Tensoring the finite-module
presentation and the injection of the parameter polynomial algebra
with a field extension preserves finiteness and injectivity. Hence
the same forms work after extending scalars.

For the unconditional bound, use the homogeneous-feasibility PSPACE
procedure cited in GCT~V Theorem~2.10. Enumerate bounded lists with
polynomial-space counters and test them using that procedure,
increasing $s$ until a successful list appears. All counters and
oracle workspaces are polynomial in $N$, giving a functional PSPACE
construction, consistent with GCT~V Theorem~3.4.
\end{proof}

The zero algebra is an excluded trivial case. The coordinate list
condition rules out encoding exponentially many unused free variables
by a binary dimension alone. Standard grading is needed for the
linear homogeneous output asserted here. The oracle bound does not
give an ordinary polynomial-time algorithm.

\subsection{CX-K06: the exact EXPH transfers for succinct explicit varieties}
\label{cx:k06}

\paragraph{Statement.}
Under $H_M$, GCT~V Theorem~5.12 gives construction and verification of
a strict succinct small system of parameters for an explicit family
$W_n$ as in Definition~5.1 in the functional exponential hierarchy;
the determinant version is Theorem~4.10. For an optimal-cardinality
h.s.o.p., the EXPH assertion in Theorem~5.5 retains the additional
hypothesis that defining equations can be computed in
$2^{\operatorname{poly}(n)}$ time; its determinant counterpart is
Theorem~4.1.

\begin{proof}
As in CX-K02, $H_M$ supplies $H_D$. Every application of Koiran in
GCT~V Theorem~2.8 concerns a finite integer or rational circuit system,
possibly of exponential encoding length. Its arithmetic premise is
therefore discharged by this same $H_D$.

For clarity about equation access, the proof of Theorem~2.8 itself
constructs an existential polynomial system of exponential total
encoding length and exponentially many variables for a candidate
border-circuit hitting set. Infeasibility of that system is the
hitting-set predicate. The construction uses
B\"urgisser \cite[Theorem~5.7]{Burgisser04}, as cited there; it does
not assume an exponentially small defining ideal for $W_n$.
The systems are constructed in deterministic exponential time.
Their good-candidate predicates have polynomial time and quantifier
bounds in the expanded length, with a constant number of
alternations. In terms of $n$, those bounds become exponential,
still with constantly many alternations. Searching a polynomial-length
description of a hitting set, followed by the transformation in
GCT~V Theorem~5.11(a), therefore stays in the functional exponential
hierarchy. The same predicates give the verification assertion in
Theorem~5.12. This proves the strict s.s.o.p. transfer and preserves
the output conventions of the cited theorem.

For the optimal-cardinality h.s.o.p. assertion, use GCT~V
Theorem~5.5 exactly as stated. Its proof invokes the argument of
Theorem~4.1 and hence Theorem~3.4. Their arithmetic premise is again
$H_D$. The extra equation algorithm supplies an expanded input of
length $L\le 2^{\operatorname{poly}(n)}$. Polynomial time and
polynomial quantifier length in $L$ become exponential time and
exponential quantifier length in $n$, while the number of alternations
stays constant. In the standard-graded case this substitution is
directly the algorithm in CX-K05. For differing positive weights,
we invoke the stated graded theorem with its representation
conventions; arbitrary linear combinations of generators of different
weights are not homogeneous.

Without equation access of the required time bound, this proof has
no first step of the stated size. The unconditional elimination
method discussed in GCT~V may produce double-exponential-length
equations. Supplying $H_D$ does not change that representation cost.
The extra assumption in Theorem~5.5 therefore remains part of this
specific EXPH conclusion.
\end{proof}

We use ``exponential hierarchy'' for the union of constant levels
defined by alternating exponential time with exponentially bounded
quantifiers; a functional version permits the corresponding oracle
construction. The argument proves membership at a fixed finite level
under the cited source conventions. It does not identify general
EXPSPACE with EXPH or claim that all normalizations require the
equation-access hypothesis.

\subsection{CX-K07: the existing unconditional PSPACE construction}
\label{cx:k07}

\paragraph{Statement.}
Polynomial-size succinct separating normalization systems for explicit
varieties already have unconditional polynomial-space constructions,
independent of $H_A$ and $H_M$. Thus their EXPH construction does not
become unconditional for the first time through either stipulated
2026 premise. Optimal-cardinality h.s.o.p.'s and verification of
arbitrary candidate systems remain different statements.

\begin{proof}
Forbes--Shpilka \cite[Theorem~7.1, preprint p.~21]{FS18} gives a
polynomial-space construction of a polynomial-size rational hitting
set, of polynomial bit complexity, for the closure of small
bounded-degree algebraic circuits over $\mathbb R$ or $\mathbb C$.
Its parameters are supplied in unary. Apply it to the polynomially
bounded parameters of the explicit family and feed the resulting set
to GCT~V Theorem~5.11(a). Polynomial-space construction followed by
the polynomial-time output transformation remains polynomial-space.

Guo--Saxena--Sinhababu \cite[Theorems~1.5 and~1.6, p.~6]{GSS19}
explicitly record the determinant and general explicit-variety
normalization consequences, and provide the corresponding finite-field
treatment. These results assume no GRH and predate the new premises.
Their outputs allow polynomial slack in the number of parameters;
they do not assert that the cardinality equals the dimension.
\end{proof}

This is a direct composition of cited algorithms and a correction to
the scope of the historical comparison. It is not a new discovery of
the PSPACE algorithm. Guo--Saxena--Sinhababu uses the term e.s.o.p.
while placing its construction in PSPACE. To avoid conflict with
GCT~V's polynomial-time convention for that term, the conclusion here
is stated as a polynomial-space construction of a polynomial-size
succinct separating system. Neither that theorem nor this composition
is an assertion about every arbitrary candidate-verification input.

\subsection{Remaining gates and verification status}

The full equation-given h.s.o.p. conclusion in
$\mathrm{FP}^{\Sigma_2^{\mathrm P}}$ remains conditional on $H_M$
in this run. From $H_A$ alone, CX-K03 establishes the explicitly
bounded abelian-witness promise version, and CX-K04 refutes the
unrestricted abelian-witness shortcut. No logical nonimplication
between $H_A$ and general HN membership is claimed. The succinct
optimal-cardinality h.s.o.p. EXPH transfer retains its equation-access
gate; the small separating-system construction already has the
unconditional PSPACE bound.

All proofs CX-K01--CX-K07 were read end to end after drafting by
\texttt{/root/koiran\_gct}, GPT-6 Astra Pro, on 20 November 2026 UTC;
the inherited effort setting was not exposed. The independent agents
\texttt{/root/eth\_consequences} and \texttt{/root} each read the
complete proofs. Their checks found no remaining substantive gap after
the standard-grading and output-scope qualifications were incorporated.
Koiran's protocol and erratum, the relevant GCT statements and displayed
proofs, and the later theorem statements were inspected. The full
historical analytic and border-hitting-set proofs are carried as cited
background, not claimed to have been independently reproved. Detailed
check events and source locators accompany the ledger.
